\section{Joint collision blocks under a shrinking endpoint window}
\label{sec:multiblock-bridge-windows}
The comparison of two conditioned initial measures does not identify
their common limit. We now retain finitely many path frequencies while
integrating the endpoint frequency. The estimates are uniform in the
positions of the block boundaries, including coincident boundaries.
This uniformity will be used for conditional tightness, not only for
finite-dimensional convergence.

Write $S_mh=S_{m,R}h_R$, $e_0=3/280$, and
$\mathfrak e_m=m^{-e_0}\sqrt{\log(2+m)}$, as in
Section~\ref{sec:stationary-collision-band}. All collision expectations
in this section are under $\sigma\nu$, where
\begin{equation}\label{eq:bridge-initial-class}
 \sigma\ge0,\qquad \int\sigma\dd\nu=1,\qquad
 \|\sigma\|_\infty+\|\sigma\|_{\mathrm{BV}}\le B.
\end{equation}
The constants may depend on the fixed number of blocks, but not on
any of their lengths.

\subsection{A uniform joint Fourier estimate}
Fix an integer $q\ge1$ and integers
$0=j_0\le j_1\le\cdots\le j_q=m$. Put
$m_l=j_l-j_{l-1}$, $d_l=m_l/m$, and
$H_l=(S_{j_l}h-S_{j_{l-1}}h)/\sqrt m$.
For $u_1,\ldots,u_q\in\R^4$ define
\begin{align}
 \mathcal C_m(v;\mathbf u)
 &=\int\sigma\exp\left(i\sum_{l=1}^q(v+u_l)\cdot H_l\right)\dd\nu,
 \notag\\
 \mathcal G_m(v;\mathbf u)
 &=\exp\left(-\frac12\sum_{l=1}^q
               d_l(v+u_l)^{\mathsf T}\Gamma_R(v+u_l)\right).
 \label{eq:multiblock-characteristic}
\end{align}
An interval of length zero contributes zero to both expressions.

\begin{theorem}[The endpoint band with joint block frequencies]
\label{thm:multiblock-collision-band}
For fixed $q,B,C_0<\infty$, uniformly in $R$, the subdivision, and
$\max_l|u_l|\le m^{1/200}$,
\begin{equation}\label{eq:multiblock-collision-band}
 \int_{|v|\le C_0m^{9/100}}
   |\mathcal C_m(v;\mathbf u)-\mathcal G_m(v;\mathbf u)|\dd v
          \le C_{q,B,C_0}\mathfrak e_m.
\end{equation}
Here $v$ and $u_l$ are rescaled frequencies. Each actual collision
frequency is $(v+u_l)/\sqrt m$.
\end{theorem}
\begin{proof}
On $|v|\le2m^{1/200}$ take $\delta=m^{-1/14}/4$ and smooth
$\sigma$ and $h_R$. Put $z_l=(v+u_l)/\sqrt m$. In chronological
order the smoothed pairing is
\[
 \ell\left(L_{z_q}^{m_q}\cdots L_{z_1}^{m_1}
                         (\mathsf S_\delta\sigma)\nu\right),
 \qquad L_z=\mathcal L_{R,\delta,z}.
\]
Use the decomposition in Lemma~\ref{lem:smooth-collision-expansion}.
For a zero power its complementary part means $I-\Pi(z)$.
The complementary bound $C\rho^l$ and its annihilation of $\Pi(z)$
then hold also at $l=0$.

We explain why a short block is harmless. In the expansion of this
fixed product, the term with only complementary factors is at most
$C_{q,B}\delta^{-2}\rho^m$. Every other nonprincipal term has two
adjacent factors of different types. At such a boundary, for example,
\[
 N(z)^{l}\Pi(w)=N(z)^{l}(\Pi(w)-\Pi(z)).
\]
The reversed product is handled in the same way. Analytic projection
bounds give $\|\Pi(w)-\Pi(z)\|\le C\delta^{-2}|w-z|$ on this
ball. All remaining powers are uniformly bounded: their eigenvalue
factors have modulus at most
$\exp(Cm\delta^{-6}\max_l|z_l|^3)\le C$, and the complementary
powers have their uniform bounds. Thus the sum of the mixed terms,
including the initial density norm, is at most
$C_{q,B}\delta^{-4}(|v|+U)/\sqrt m$, where $U=\max_l|u_l|$.
This argument uses no positive lower bound for $m_l/m$.

The product of principal factors has amplitude
$1+O_{q,B}(\delta^{-4}(|v|+U)/\sqrt m)$. The eigenvalue product
has the Gaussian logarithm with $\Gamma_{R,\delta}$, up to a remainder
bounded by $C_q\delta^{-6}m^{-1/2}(|v|+U)^3$.
Changing the density costs $C_B\delta$. For the observable change,
the small-mass variance on each deterministic interval and Minkowski's
inequality give
\[
 \int\left|\sum_l z_l\cdot
       (S_{j_l}-S_{j_{l-1}})(h_R-h_{R,\delta})\right|
                       (\mathsf S_\delta\sigma)\dd\nu
 \le C_{q,B}(|v|+U)\sqrt{\delta(1+|\log\delta|)}.
\]
Changing the covariance costs
$C_q\delta(1+|\log\delta|)(|v|+U)^2$. Since
$U\le m^{1/200}$, integration of these bounds has exactly the six
power budgets in the proof of Theorem~\ref{thm:wide-collision-band}.
The slowest is $1/28-5/200=3/280$. This proves the central estimate,
uniformly in the permitted growing block frequencies.

On $2m^{1/200}\le|v|\le C_0m^{9/100}$ use the fixed orders and
scales
\[
 P=40,\quad Q=29,\quad \delta=m^{-19/100}/4,
                   \quad\epsilon=m^{-50}/4.
\]
Here $|v+u_l|\ge|v|/2$ and $|v+u_l|\le3|v|/2$. We give the
multiblock version of the argument in
Theorem~\ref{thm:high-order-damped-unsmoothing} to account for the
additional frequencies. Let $r_\delta=h_R-h_{R,\delta}$ and
$Z=\max_l|z_l|$. For $Z>0$, write the residual phase as
\[
 X=\sum_l z_l\cdot(S_{j_l}-S_{j_{l-1}})r_\delta.
\]
The all-small-mass moment bound, on each block, followed by
Minkowski's inequality gives
\[
 \int|X|^{2P}\dd\nu
       \le C_{q,P}Z^{2P}\sum_{b=1}^P(m\delta)^b
                                (1+|\log\delta|)^{2P-b}.
\]
Domination gives the same bound with the smoothed initial density.
Expand $e^{iX}$ through degree $79$. Smooth its residual factors
and the initial density at scale $\epsilon$. Each corrected monomial
is a chronological word with at most $80$ smooth multipliers and the
$q-1$ additional block boundaries. The power lengths in block $l$
sum to $m_l$, so its damping is
\[
 \exp\left(-c\sum_l m_l|z_l|^2\right)
                                      \le e^{-c'|v|^2}.
\]
The number of factors is bounded in terms of $q$ and the fixed order,
not in terms of $m$. The fine insertion, fine residual, and residual
moment errors are consequently bounded exactly as in
\eqref{eq:high-order-damped-unsmoothing}, with $|z|$ replaced by $Z$
and a changed fixed constant. The two analytic margins are $3/100$
and $2/25$. The integrated fine errors have decay exponents
$1241/25$ and $303/100$. The largest residual-moment power is
$-1/25$, attained at $b=40$; every smaller block number improves it
by $81/100$. Thus the margin $41/1400$ pays all fixed logarithms.
The Gaussian in \eqref{eq:multiblock-characteristic} has the same
annular damping. Combining the two regions proves the theorem for
large $m$; enlarging the constant covers the finitely many remaining $m$.
\end{proof}

\subsection{A shrinking endpoint with oscillatory path weights}
Let $P_m:\R^4\to\R^3$ be the first three rows of a matrix
$\mathsf A_m$ satisfying
$\|\mathsf A_m\|+\|\mathsf A_m^{-1}\|\le C_0$. Put
\[
 C_m=P_m\Gamma_RP_m^{\mathsf T},\qquad
 Y_m=P_mS_mh/\sqrt m,\qquad
 E_m=\{Y_m\in B_3(\xi,\mathbf h)\},
\]
where $|\xi|\le A$ and
$c_0m^{-2/25}\le h_i\le C_0m^{-2/25}$. These assumptions make
$C_m$ uniformly positive definite. Write $V_3=\prod_i2h_i$ and
\[
 \varepsilon_m=m^{-23/2800}\sqrt{\log(2+m)}.
\]
For $\theta_1,\ldots,\theta_q\in\R^3$ define
\begin{align}
 \mathcal B_m(\boldsymbol\theta;y)
 =\exp\biggl(&i\sum_l d_l\theta_l\cdot y
 -\frac12\sum_{l,r}(d_l\mathbf1_{l=r}-d_ld_r)
                         \theta_l^{\mathsf T}C_m\theta_r\biggr).
 \label{eq:bridge-increment-characteristic}
\end{align}
This is the characteristic function of Gaussian increments of lengths
$d_l$, conditioned on their sum being $y$.

\begin{proposition}[Uniform conditional block characteristic functions]
\label{prop:window-conditional-characteristic}
There is $c(C_0)>0$ such that
$\max_l|\theta_l|\le c(C_0)m^{1/200}$ implies
\begin{align}
 &\left|\mathbb E_{\sigma\nu}
  \left[\left.\exp\left(i\sum_l\theta_l\cdot P_mH_l\right)
                          \right|E_m\right]
           -\mathcal B_m(\boldsymbol\theta;\xi)\right|
       \le C_{q,A,B,c_0,C_0}\varepsilon_m .
 \label{eq:window-conditional-characteristic}
\end{align}
The bound is uniform over all subdivisions, including zero-length
blocks. It concerns the original endpoint indicator, with either
choice of inclusion of its discrete endpoints.
\end{proposition}
\begin{proof}
First retain the Gaussian integral over the endpoint box rather than
its center. Augment $Y_m$ by the fourth row of $\mathsf A_m$ and
truncate that fourth coordinate to $[-T,T]$, $T=m^{1/400}$.
Use the four-dimensional endpoint envelopes of
\eqref{eq:box-lower-envelope} at bandwidth $b=c_1(C_0)m^{9/100}$.
After the change of frequency $v\mapsto\mathsf A_m^{\mathsf T}v$,
the oscillatory comparison is
Theorem~\ref{thm:multiblock-collision-band}, with
$u_l=P_m^{\mathsf T}\theta_l$. Its error for either envelope is
$C T V_3\mathfrak e_m$.

Positivity cannot sandwich an oscillatory numerator. Instead let
$W=\exp(i\sum_l\theta_l\cdot P_mH_l)$, so $|W|=1$, and use
\[
 |\mathbb E[W(\one_{B_3\times[-T,T]}-U)]|
                  \le\mathbb E(U-\one_{B_3\times[-T,T]})
                  \le\mathbb E(U-L).
\]
The latter is bounded by the unweighted Fourier comparison plus the
Gaussian envelope excess. The Gaussian weighted endpoint measure
also has total variation bounded by its unweighted Gaussian marginal,
because its conditional characteristic function has modulus at most
one. Thus the same replacement is valid on the Gaussian side.
Integrating the fourth Gaussian coordinate against its decay, as in
\eqref{eq:stationary-projection-budget}, gives an error
$C V_3[\sum_i(bh_i)^{-1}+b^{-1}]$, without a factor $T$.
The $128$th collision moment bounds the omitted actual tail by
$CT^{-128}$; $|W|=1$ makes this estimate unchanged. The Gaussian
tail is at most $C V_3e^{-cT^2}$.

It follows that the numerator differs from
\[
 \int_{B_3(\xi,\mathbf h)}g_{C_m}(y)
                         \mathcal B_m(\boldsymbol\theta;y)\dd y
\]
by at most $C V_3\varepsilon_m$. The three exponent margins are
$23/2800$, $1/100$ and $2/25$, exactly as in the unweighted projection.
The Gaussian formula follows by subtracting $d_l$ times the total
from independent centered increments of covariance $d_l C_m$.
These differences have covariance
$(d_l\mathbf1_{l=r}-d_ld_r)C_m$ and are independent of the total.

Theorem~\ref{thm:projected-collision-windows} gives
$\mathbb P_{\sigma\nu}(E_m)=V_3g_{C_m}(\xi)(1+O(\varepsilon_m))$,
including a lower bound $c_A V_3$. On the box the change in
\eqref{eq:bridge-increment-characteristic} from its value at $\xi$
is at most $C m^{-2/25}\max_l|\theta_l|$; the nonoscillatory
Gaussian density changes by $O_A(m^{-2/25})$. The first error is
$O(m^{-3/40})$, hence smaller than $\varepsilon_m$ even for the
permitted growing frequencies. Divide by the exact denominator.
This proves the proposition, without taking a trace of a
four-dimensional $L^1$ Fourier bound.
\end{proof}

\paragraph{Uniformity needed below.}
The number of blocks in these statements is fixed. Their boundaries
may nevertheless depend on $m$ and may approach each other or an
endpoint. The permitted growth of the test frequencies and this
boundary uniformity are both needed to prove tightness after division
by a probability of order $m^{-6/25}$. Finite-dimensional convergence
alone would not justify such a division in a path-space assertion.
